% !TeX root = ../动手学强化学习笔记.tex
% 由主文档通过 \input 载入；请编译主文档。
\setcounter{chapter}{18} % 保留原书章号。
\chapter{探索、奖励设计与约束决策}
\label{ch:rl-exploration}

探索决定能够获得什么数据，奖励决定算法用什么标准比较行为，约束则规定哪些策略可接受。三者相互影响，但不能互相替代：即使优化算法完全正确，缺失的探索或错误的奖励也可能导致不希望出现的行为。本章将三者作为相邻而有区别的专题：探索部分说明轨迹结构与内在奖励的作用，奖励设计部分证明势函数塑形在相应条件下保持策略排序，约束部分仅建立约束 MDP 与拉格朗日求解思路。学习时应分别判断一个改动改变了采样过程、奖励目标还是可行策略集合；约束算法的可行性与收敛理论属于进一步学习的范围。

\section{从局部探索到长时程探索}

$\varepsilon$-贪心每步独立随机化，在需要连续采取特定动作的任务中可能效率很低。若只有连续 $L$ 次向右才有奖励，独立公平左右选择一次成功的概率为 $2^{-L}$。在整段尝试中保持一个探索方向，则改变了到达远处状态的概率。这说明探索需要关注轨迹结构，而不仅是单步动作熵。

访问计数奖励如 $\beta/\sqrt{N(s,a)+1}$ 偏向少见状态动作，其中 $\beta>0$ 决定尺度；乐观价值估计偏向后果尚不确定但可能有利的动作；信息收益关注减少关于环境的后验不确定性。三者衡量的对象不同，不能把奖励高、访问少和信息多视为同一件事。高维连续状态还需要定义相似性，否则每个观测都可被判为“全新”。

\section{内在奖励与目标偏移}

内在奖励（intrinsic reward）是为支持探索而构造的辅助信号。若训练奖励为 $r^{\rm train}=r^{\rm task}+\beta r^{\rm int}$，算法实际优化的是组合目标。评价任务回报时应使用 $r^{\rm task}$，否则可能把持续访问新奇区域误认为完成了任务。

预测误差是一种新奇性代理，但随机噪声源即使被充分了解也可能始终难以预测。例如不可控制的随机屏幕不断变化，像素预测误差长期较高，却不一定提供有助于决策的信息。应区分无法消除的随机性与能够通过学习减少的知识不确定性。

\section{势函数奖励塑形}

\begin{theorem}[折扣任务中的势函数塑形]
设 $0\leq\gamma<1$，原奖励有界，$\Phi:\mathcal S\to\mathbb R$ 有界。将奖励改为 $R'_{t+1}=R_{t+1}+\gamma\Phi(S_{t+1})-\Phi(S_t)$，则对任意策略和固定初始状态 $s$，新价值为 $v'_\pi(s)=v_\pi(s)-\Phi(s)$，因此逐状态的策略排序不变。
\end{theorem}
\begin{proof}
前 $T$ 步新增折扣和为
\[
\sum_{t=0}^{T-1}\gamma^t\bigl(\gamma\Phi(S_{t+1})-\Phi(S_t)\bigr)
=-\Phi(S_0)+\gamma^T\Phi(S_T).
\]
末项因有界性趋于零。新增回报只依赖起点而不依赖策略，取期望即得结论。
\end{proof}

这称为势函数奖励塑形（potential-based reward shaping）。例如令势函数随接近目标而增大，可以向中间步骤提供较密反馈，同时在上述条件下保持原最优策略。有限回合若在终止时直接截断求和，则仍有 $\gamma^T\Phi(S_T)$，通常应把终端势设为零。无折扣随机终止情形也要单独控制终端项和可积性。

\begin{example}[忽略边界会改变最优动作]
一步任务中两动作原奖励均为零，分别到达两个终止状态，其势为 $0$ 和 $1$。若 $\gamma=1$ 且起点势为零，塑形后奖励分别为 $0,1$，原本并列的动作不再并列。错误来自未处理终端势，而不是望远镜求和公式失效。
\end{example}

\section{约束 MDP 与目标取舍}

约束马尔可夫决策过程（constrained MDP）同时指定奖励目标 $J(\pi)$ 和成本信号 $C_{t+1}\geq0$。例如定义折扣期望成本 $D(\pi)=\mathbb E_\pi\sum_t\gamma^tC_{t+1}$，求解 $\max_\pi J(\pi)$，满足 $D(\pi)\leq d_0$。拉格朗日函数为 $\mathcal L(\pi,\lambda)=J(\pi)-\lambda(D(\pi)-d_0)$，$\lambda\geq0$。可交替提高策略的惩罚后目标，并按 $\lambda\leftarrow[\lambda+\beta(\widehat D-d_0)]_+$ 调整乘子。

该迭代只是求解思路；可行性、强对偶与收敛需额外条件，非凸神经网络训练不能自动获得这些性质。固定一个有限罚系数也不等于严格约束。期望成本小仍可能允许少量轨迹具有很大成本；每步硬约束、失败概率限制和尾部风险度量分别对应不同数学要求，应根据任务选择。

\section{习题与实验}
\begin{enumerate}
\item 对长度 $L$ 的稀疏奖励链比较独立随机动作和整段保持方向的成功概率，写明重置和失败规则。
\item 证明从任意时刻 $t$ 起的塑形回报差为 $-\Phi(S_t)$，并解释对动作价值和优势的影响。
\item 设计带危险区域的网格，分别测试辅助探索奖励、势函数塑形和成本约束，报告原任务回报与成本，避免把三种机制混为同一个改动。
\end{enumerate}





